\documentclass[11pt]{article}
\usepackage[T1]{fontenc}
\usepackage{textcomp}
\usepackage[margin=0.75in]{geometry}
\usepackage{amsmath,amssymb,amsthm}
\usepackage{array,booktabs}
\usepackage{hyperref}
\setlength{\parindent}{0pt}
\newtheorem{theorem}{Theorem}
\theoremstyle{definition}
\newtheorem{definition}{Definition}
\title{Flow Proof Artifact: prop-11-gnomon-applications}
\author{Generated by \texttt{flow doc proof}}
\date{Flow Proof Book}
\begin{document}
\maketitle
To apply a gnomon equal to a given rectangle in a parallelogram about a given straight line.\\[0.5em]
\textbf{Source.} euclid — Elements, Book II, Proposition 11\\[0.5em]
\bigskip
\noindent{\large \textbf{Derived fact 1.} \textit{To apply a gnomon equal to a given rectangle in a parallelogram about a given straight line} \hfill the Euclidean plane \cdot Euclid Book II \cdot Proposition 11: a gnomon equals the rectangle by the segments}\par
\smallskip\noindent\textit{Coordinate: the Euclidean plane \cdot Euclid Book II \cdot Proposition 11: a gnomon equals the rectangle by the segments}\par
\smallskip\noindent\textit{Source: euclid — Elements, Book II, Proposition 11}\par
\smallskip\noindent\textit{Built on: proposition 34: complements of parallelograms about the diameter are equal, for Book I of the Elements in on the Euclidean plane, proposition 1: the rectangle by two lines equals the sum over segments, for Euclid Book II on the Euclidean plane}\par
\medskip
\noindent\textbf{Goal.} To apply a gnomon equal to a given rectangle in a parallelogram about a given straight line\par
\noindent$\displaystyle \text{gnomon equals given rectangle}$\par
\medskip
\noindent
\begin{tabular}{@{} >{\raggedright\arraybackslash}p{0.06\textwidth} >{\raggedright\arraybackslash}p{0.47\textwidth} >{\raggedleft\arraybackslash}p{0.41\textwidth} @{}}
\textbf{\#} & \textbf{Proof} & \textbf{Mathematics} \\
\midrule
\textbf{1.} & We prove that proposition 11: a gnomon equals the rectangle by the segments for Euclid Book II the Euclidean plane. & \\[0.45em]
\textbf{2.} & Let target = given\_rectangle. & \\[0.45em]
\textbf{3.} & Let base = given\_straight\_line. & \\[0.45em]
\textbf{4.} & Let gnomon = gnomon\_about\_diameter. & \\[0.45em]
\textbf{5.} & From step 2, step 3, and step 4, we can deduce that gnomon equals target. & $\displaystyle gnomon = target$ \\[0.45em]
\textbf{6.} & We invoke the derived fact governing Book I of the Elements in the Euclidean plane: proposition 34: complements of parallelograms about the diameter are equal, for Book I of the Elements in on the Euclidean plane. & \\[0.45em]
\textbf{7.} & We invoke the derived fact governing Euclid Book II the Euclidean plane: proposition 1: the rectangle by two lines equals the sum over segments, for Euclid Book II on the Euclidean plane. & \\[0.45em]
\textbf{8.} & From step 6 and step 7, this implies gnomon equals given rectangle. Hence proven. & $\displaystyle \text{gnomon equals given rectangle}$ \\[0.45em]
\bottomrule
\end{tabular}
\label{thm:Geometry-Euclid Book II-proposition_11_a_gnomon_equals_the_rectangle_by_the_segments}
\par\medskip\hrule
\end{document}
